% !TeX root = ../../Book4.tex
% 纲要标题：### E.1 覆盖、景与集合值层
% 纲要细目（写作范围）：
% #### E.1.1 筛与 Grothendieck 拓扑
% #### E.1.2 层的相容族与唯一拼接
% #### E.1.3 开集覆盖与群作用的例子

\section{覆盖、景与集合值层}
\label{b4:sec:E01}

在拓扑空间上，局部数据的相容性由开集的交来检验。在域上，添入一个可分根也可以成为局部变换，却不是原空间中的开子集。要同时描述这两种情形，须把“哪些态射足以覆盖一个对象”作为数据。筛把覆盖连同它的全部限制放在一起，使定义不要求范畴中处处存在纤维积。

\subsection{筛与 Grothendieck 拓扑}
\label{b4:subsec:E01:01}

\begin{definition}\label{b4:def:sieve}
设 \(\mathcal C\) 为小范畴，\(U\in\mathcal C\)。一族以 \(U\) 为目标的态射 \(S\)，若满足：\(f:V\to U\) 属于 \(S\) 时，每个可复合的 \(g:W\to V\) 都有 \(f\circ g\in S\)，则称 \(S\) 为 \(U\) 上的\term[shai]{筛}[sieve]。等价地，它是可表预层 \(h_U=\operatorname{Hom}_{\mathcal C}(-,U)\) 的子预层。对 \(f:V\to U\)，定义
\[
 f^*S=\{\,g:W\to V\mid f\circ g\in S\,\}.
\]
这称为筛沿 \(f\) 的拉回。
\end{definition}
\symidx{\(f^*S\)：筛沿态射的拉回}

由一族箭头 \((u_i:U_i\to U)\) 生成的筛，包括所有能经过某个 \(u_i\) 分解的箭头。因而它记录的不只是这组箭头，还包括由它们继续限制得到的局部数据。最大筛为整个 \(h_U\)，其中包含 \(1_U\)。

\ProseParagraphBreak
例如，把链 \(0<1<2\) 看成偏序范畴。在对象 \(2\) 上，由 \(1\to2\) 生成的筛还必须包含 \(0\to2\)，因为后者经过 \(0\to1\to2\) 分解；它不包含 \(1_2\)，因而不是最大筛。沿 \(1\to2\) 拉回时，\(1_1\) 与 \(0\to1\) 都在拉回筛中，所得却是对象 \(1\) 上的最大筛。筛的封闭条件把一张箭头及它的全部进一步限制一起保留。

\begin{definition}\label{b4:def:grothendieck-site}
设 \(\mathcal C\) 为小范畴。对每个 \(U\in\mathcal C\)，指定一族筛 \(J(U)\)，并满足以下条件：
\begin{enumerate}
\item 最大筛属于 \(J(U)\)。
\item 若 \(S\in J(U)\)、\(f:V\to U\)，则 \(f^*S\in J(V)\)。
\item 若 \(S\in J(U)\)，且 \(R\) 为 \(U\) 上的筛，对每个 \(f:V\to U\) 属于 \(S\) 都有 \(f^*R\in J(V)\)，则 \(R\in J(U)\)。
\end{enumerate}
称 \(J\) 为 \(\mathcal C\) 上的\term[Grothendiecktuopu]{Grothendieck 拓扑}[Grothendieck topology]，\((\mathcal C,J)\) 为\term[jing-site]{景}[site]，\(J(U)\) 中的筛为覆盖筛。
\end{definition}
\symidx{\(J(U)\)：Grothendieck 拓扑指定的覆盖筛族}

第二条说覆盖在改变观察对象后仍是覆盖；第三条说先覆盖、再在每块上覆盖，能够组成覆盖。这种“拓扑”指定的是态射的覆盖关系，不是给对象集合添上一族普通开集。定义及一般构造见 Mac Lane 与 Moerdijk\cite[Chapter III, Sections 1--2]{MacLaneMoerdijk1992Sheaves}。

\ProseParagraphBreak
三条公理还推出覆盖筛对包含关系向上封闭：若 \(S\in J(U)\) 且 \(S\subseteq R\subseteq h_U\)，则 \(R\in J(U)\)。事实上，对每个 \(f:V\to U\) 属于 \(S\)，也有 \(f\in R\)，故 \(1_V\in f^*R\)。筛中一旦含恒等箭头，就包含以该对象为目标的全部箭头，所以 \(f^*R=h_V\) 是覆盖筛。第三条公理随即给出 \(R\in J(U)\)。增加已有覆盖的箭头及其限制，不会破坏覆盖性质。

\subsection{层的相容族与唯一拼接}
\label{b4:subsec:E01:02}

\begin{definition}\label{b4:def:site-set-sheaf}
设 \((\mathcal C,J)\) 为小景，\(F:\mathcal C^{\mathrm{op}}\to\mathbf{Set}\) 为预层。若对每个 \(U\) 及每个 \(S\in J(U)\)，限制映射
\[
 F(U)\longrightarrow\operatorname{Nat}(S,F),\qquad
 x\longmapsto\bigl(f\longmapsto F(f)(x)\bigr)
\]
为双射，则称 \(F\) 为该景上的\term[jihezhiceng]{集合值层}[sheaf of sets]。这些层及其自然变换组成的范畴记为 \(\operatorname{Sh}_{\mathbf{Set}}(\mathcal C,J)\)。
\end{definition}
\symidx{\(\operatorname{Sh}_{\mathbf{Set}}(\mathcal C,J)\)：景上的集合值层范畴}

右边的一项自然变换，就是对每个 \(f:V\to U\) 属于 \(S\) 指定 \(x_f\in F(V)\)，且有
\[
 x_{f\circ g}=F(g)(x_f).
\]
双射的单射性表示局部限制唯一决定整体，满射性表示相容的局部数据能够拼接。\(\mathcal C\) 小保证这里的自然变换构成集合。

\ProseParagraphBreak
若 \(S\) 由覆盖箭头 \(u_i:U_i\to U\) 生成，且各纤维积 \(U_i\times_UU_j\) 存在，那么筛上的相容族等价于一族 \(x_i\in F(U_i)\)，满足
\[
 F(p_1)(x_i)=F(p_2)(x_j),
\]
其中 \(p_1,p_2\) 是两张投影，等式在 \(F(U_i\times_UU_j)\) 中成立。自然变换沿 \(u_ip_1=u_jp_2\) 限制，便给出这个等式。反过来，对 \(f:V\to U\) 属于 \(S\)，选分解 \(f=u_i g\)，令 \(x_f=F(g)(x_i)\)。若还可写为 \(f=u_j h\)，纤维积的泛性质给出 \(V\to U_i\times_UU_j\)，把上式沿它限制就得到 \(F(g)(x_i)=F(h)(x_j)\)，所以定义与分解的选择无关；继续前复合也保持这一定义。这里的条件包括 \(i=j\)：同一覆盖箭头的不同分解也必须相容。

\begin{proposition}\label{b4:prop:open-site-sheaves}
设 \(X\) 为拓扑空间，以开集及包含为范畴 \(\mathcal O(X)\)。规定 \(U\) 上的筛覆盖，当且仅当其中所有开集的并为 \(U\)。这给出 Grothendieck 拓扑；集合值层条件等价于通常开覆盖上的相容截面唯一拼接条件。
\end{proposition}
\symidx{\(\mathcal O(X)\)：拓扑空间的开集范畴}
\begin{proof}
筛是在包含关系下向下封闭的开集族。最大筛显然覆盖。沿 \(V\subseteq U\) 拉回筛时，其中的开集覆盖 \(V\)，因为可用 \(V\cap W\) 限制原覆盖中的 \(W\)。若 \(S\) 覆盖 \(U\)，且 \(R\) 在每个 \(W\in S\) 上的拉回覆盖 \(W\)，则 \(R\) 中开集的并包含每个 \(W\)，故等于 \(U\)。这核验三条公理。

开集包含的纤维积为 \(U_i\cap U_j\)，所以前述相容等式就是截面在交集上的限制相同。任意覆盖筛本身也是开覆盖，故筛上的唯一拼接与开覆盖上的唯一拼接相互推出。空集由空族覆盖，还得到 \(F(\varnothing)\) 是单元素集合。
\end{proof}

第三卷第24.1--24.2节的局部化构造就是这一条件的具体实例。设 \(A\) 为交换环、\(M\) 为 \(A\) 模，\(X=\operatorname{Spec}A\)。结构层与伴随模层在基本开集上的截面分别为 \(\mathcal O_X(D(f))=A_f\) 与 \(\widetilde M(D(f))=M_f\)，限制由继续局部化给出。若 \(X=\bigcup_iD(f_i)\)，在 \(D(f_i)\cap D(f_j)=D(f_if_j)\) 上限制相同的局部分式便唯一黏合成全局截面。把这组基本开覆盖及其所有开子集放入生成的筛，就得到上一定义中的相容族；环乘法与模作用是这些集合值层上的附加结构。

\ProseParagraphBreak
这里取值于集合。\secref{b4:sec:C03}中的层取值于交换群，截面另有加法；不能仅因同样满足拼接就把两个层范畴当作同一个范畴。

\subsection{开集覆盖与群作用的例子}
\label{b4:subsec:E01:03}

若每个对象只有最大筛被规定为覆盖，则任意预层都是层，因为 Yoneda 引理给出 \(\operatorname{Nat}(h_U,F)\cong F(U)\)。这称为平凡的 Grothendieck 拓扑。若 \(G\) 为群，把它看成单对象范畴，则每个非空筛都含全部箭头：含 \(g\) 就含 \(gg^{-1}=1\)，继而含每个箭头。因此在平凡拓扑下，层范畴就是该范畴上的预层范畴，也就是右 \(G\)-集范畴；用 \(g\cdot x=x\cdot g^{-1}\) 可转为左 \(G\)-集。

\ProseParagraphBreak
这个单对象范畴只使用 \(G\) 的抽象群结构。若 \(G\) 还带拓扑，要让它作用于离散集合时连续，须另外要求每个元素的稳定子群都开。离散集合上的作用连续，当且仅当每个轨道映射 \(g\mapsto x\cdot g\) 连续，也就是每个点的原像开。其各个非空纤维正是稳定子群的陪集；其中 \(x\) 的原像就是稳定子群本身，而平移保持开集，故这些纤维全开当且仅当稳定子群开。例如，固定素数 \(p\)，\(\mathbb Z_p\) 在自身底集合上的右平移作用 \(x\cdot g=x+g\) 是一个抽象右 \(\mathbb Z_p\)-集。但每个元素的稳定子群都是 \(\{0\}\)，它不含任何 \(p^r\mathbb Z_p\)，因而不是开子群；把该底集合赋予离散拓扑以后，这个作用便不连续。后面的 Galois 层所对应的是离散连续作用，不能直接用全部抽象 \(G_K\)-集替代。

\ProseParagraphBreak
\subsecref{b4:subsec:C04:02}的有限平展对象采用联合满射的覆盖。对有限连续 \(G_K\)-集，这类覆盖允许把等变函数逐块拼接。开覆盖与平展覆盖虽然来源不同，层公理所做的事相同：只有与限制相容的数据才允许成为一个整体。后面研究的是这些集合值层范畴共同具有的范畴结构。
